In this work, we introduced SCOPE, an efficient hierarchical planning framework that uses LLM-generated subgoals only once at initialization, derived from suboptimal demonstration trajectories. Instead of repeatedly querying an LLM to adaptively produce subgoals during training, SCOPE directly extracts these imperfect subgoal sequences and uses them to pretrain a student planner, followed by RL-based refinement using a world model. Although the subgoals are not optimal due to the lack of interaction between the LLM and the environment, our experiments on TextCraft show that they still provide a decent starting point for hierarchical goal decomposition. Empirically, our method outperforms the LLM-based ADaPT system \citep{PrasadKHCSBK2024}, increasing the success rate from 0.52 to 0.56. It also significantly improves efficiency: to complete the game, our agent runs in 3.0 seconds on a single NVIDIA A10 GPU, whereas ADaPT requires 164.4 seconds using a GPT-3.5 backend accessed through the OpenAI API with ideal network conditions. Overall, these results demonstrate that even suboptimal one-time LLM guidance can be highly effective when combined with RL fine-tuning, offering a practical and computationally efficient alternative to LLM-dependent hierarchical planning methods.